\section{Noise Reduction and Signal Enhancement (Chapter 8.3)}

%%%%%%%%%%%%%%%%%%%%%%%%%%%%%%%%%%%%%%%%%%%%%%%%%%%%%%%%%%%%
\subsection{Grundproblem der Signalverbesserung}

In der Praxis wird selten ein ideales Signal gemessen.

Stattdessen erhält man:

\[
x(n)=s(n)+v(n)
\]

wobei

\begin{itemize}
\item $s(n)$ = gewünschtes Nutzsignal
\item $v(n)$ = Störsignal (Noise)
\end{itemize}

Das Ziel aller Verfahren aus Kapitel 8.3 ist:

\[
y(n)\approx s(n)
\]

also:

\begin{itemize}
\item Nutzsignal erhalten
\item Rauschen entfernen
\end{itemize}

\begin{center}
\includegraphics[width=0.8\linewidth]{images/noise_reduction_overview.png}
\end{center}

%%%%%%%%%%%%%%%%%%%%%%%%%%%%%%%%%%%%%%%%%%%%%%%%%%%%%%%%%%%%
\subsection{Prinzip der Noise Reduction}

Viele Störungen besitzen andere Frequenzanteile als das Nutzsignal.

Daher kann man geeignete Filter einsetzen.

\begin{center}
\includegraphics[width=0.8\linewidth]{images/signal_noise_spectrum.png}
\end{center}

Typische Beispiele:

\begin{table}[h]
\centering
\begin{tabular}{c|c}
Signal & Störung \\
\hline
ECG & 50/60 Hz Netzbrummen \\
Sprache & Hochfrequentes Rauschen \\
Radar & Zufälliges Messrauschen \\
Sensoren & Quantisierungsrauschen \\
\end{tabular}
\end{table}

%%%%%%%%%%%%%%%%%%%%%%%%%%%%%%%%%%%%%%%%%%%%%%%%%%%%%%%%%%%%
\subsection{Noise Reduction Ratio (NRR)}

Die wichtigste Größe des gesamten Kapitels ist die

\textbf{Noise Reduction Ratio (NRR)}.

Sie beschreibt, wie viel Rauschleistung nach dem Filter übrig bleibt.

Definition:

\[
NRR=
\frac{\sigma_y^2}
{\sigma_x^2}
\]

mit

\[
\sigma_x^2
\]

als Eingangsrauschleistung und

\[
\sigma_y^2
\]

als Ausgangsrauschleistung.

Für weißes Rauschen gilt:

\[
NRR=
\sum_{n=0}^{\infty} h^2(n)
\]

Alternativ:

\[
NRR=
\frac{1}{2\pi}
\int_{-\pi}^{\pi}
|H(\omega)|^2 d\omega
\]

\paragraph{Interpretation}

\begin{table}[h]
\centering
\begin{tabular}{c|c}
NRR & Bedeutung \\
\hline
$NRR<1$ & Rauschen wird reduziert \\
$NRR=1$ & keine Änderung \\
$NRR>1$ & Rauschen wird verstärkt \\
\end{tabular}
\end{table}

\paragraph{Prüfungsrelevant}

Bei Aufgabe 8.26 soll gezeigt werden:

\[
\sigma_y^2=
\sigma_x^2
\sum h^2(n)
\]

Daraus folgt direkt die NRR.

%%%%%%%%%%%%%%%%%%%%%%%%%%%%%%%%%%%%%%%%%%%%%%%%%%%%%%%%%%%%
\subsection{Signal-to-Noise Ratio (SNR)}

Ein weiteres wichtiges Maß ist das Signal-Rausch-Verhältnis:

\[
SNR=
\frac{E[s^2(n)]}
{E[v^2(n)]}
\]

Je größer die SNR, desto besser ist die Signalqualität.

Zusammenhang:

\[
\frac{SNR_{out}}
{SNR_{in}}
==========

\frac{1}{NRR}
\]

wenn das Nutzsignal durch den Filter unverändert bleibt.

%%%%%%%%%%%%%%%%%%%%%%%%%%%%%%%%%%%%%%%%%%%%%%%%%%%%%%%%%%%%
\subsection{Comb Filter}

Comb Filter gehören zu den wichtigsten Filtern des Kapitels.

Die Grundidee:

Das Signal wird mit einer verzögerten Version kombiniert.

\[
y(n)=x(n)+x(n-D)
\]

oder

\[
y(n)=x(n)-x(n-D)
\]

\paragraph{Transferfunktion}

\[
H(z)=1+z^{-D}
\]

bzw.

\[
H(z)=1-z^{-D}
\]

\paragraph{Impulsantwort}

\[
h(n)=\delta(n)+\delta(n-D)
\]

oder

\[
h(n)=\delta(n)-\delta(n-D)
\]

%%%%%%%%%%%%%%%%%%%%%%%%%%%%%%%%%%%%%%%%%%%%%%%%%%%%%%%%%%%%
\subsection{Warum heißt er Comb Filter?}

Setzt man

\[
z=e^{j\omega}
\]

ein, erhält man regelmäßig verteilte Maxima und Minima.

Der Frequenzgang besitzt die Form:

\begin{center}
\includegraphics[width=0.7\linewidth]{images/comb_response.png}
\end{center}

Er erinnert an die Zähne eines Kamms.

Daher:

\[
\text{Comb Filter}
\]

%%%%%%%%%%%%%%%%%%%%%%%%%%%%%%%%%%%%%%%%%%%%%%%%%%%%%%%%%%%%
\subsection{Nullstellen eines Comb Filters}

Für

\[
H(z)=1-z^{-D}
\]

ergibt sich:

\[
1-z^{-D}=0
\]

also:

\[
z^D=1
\]

Die Nullstellen liegen bei:

\[
z_k=e^{j2\pi k/D}
\]

\[
k=0,\dots,D-1
\]

gleichmäßig auf dem Einheitskreis.

\begin{center}
\includegraphics[width=0.5\linewidth]{images/comb_zeros.png}
\end{center}

%%%%%%%%%%%%%%%%%%%%%%%%%%%%%%%%%%%%%%%%%%%%%%%%%%%%%%%%%%%%
\subsection{NRR eines Comb Filters}

Für

\[
H(z)=1-z^{-D}
\]

ist

\[
h(n)=\delta(n)-\delta(n-D)
\]

Somit:

\[
NRR
===

# 1^2+(-1)^2

2
\]

Der Filter reduziert also kein weißes Rauschen.

Dies ist ein häufiges Prüfungsresultat.

%%%%%%%%%%%%%%%%%%%%%%%%%%%%%%%%%%%%%%%%%%%%%%%%%%%%%%%%%%%%
\subsection{Lowpass Noise Reduction Filter}

Ein einfacher Rauschfilter lautet:

\[
H(z)=
\frac{b}
{1-az^{-1}}
\]

mit

\[
0<a<1
\]

Impulsantwort:

\[
h(n)=ba^nu(n)
\]

%%%%%%%%%%%%%%%%%%%%%%%%%%%%%%%%%%%%%%%%%%%%%%%%%%%%%%%%%%%%
\subsection{Wahl des Parameters b}

Soll ein Gleichsignal unverändert bleiben:

\[
H(1)=1
\]

Daraus folgt:

\[
b=1-a
\]

Dies ist das Ergebnis von Aufgabe 8.31.

%%%%%%%%%%%%%%%%%%%%%%%%%%%%%%%%%%%%%%%%%%%%%%%%%%%%%%%%%%%%
\subsection{NRR des Lowpass-Filters}

Einsetzen der Impulsantwort ergibt:

\[
NRR
===

\sum_{n=0}^{\infty}
b^2 a^{2n}
\]

Geometrische Reihe:

\[
NRR=
\frac{b^2}
{1-a^2}
\]

Mit

\[
b=1-a
\]

folgt:

\[
NRR=
\frac{1-a}
{1+a}
\]

\paragraph{Interpretation}

\begin{table}[h]
\centering
\begin{tabular}{c|c}
a & NRR \\
\hline
0.2 & 0.67 \\
0.5 & 0.33 \\
0.9 & 0.053 \\
0.99 & 0.005 \\
\end{tabular}
\end{table}

Großes (a) bedeutet:

\begin{itemize}
\item bessere Rauschunterdrückung
\item langsamere Reaktion
\end{itemize}

%%%%%%%%%%%%%%%%%%%%%%%%%%%%%%%%%%%%%%%%%%%%%%%%%%%%%%%%%%%%
\subsection{Notch Filter}

Ein Notch Filter entfernt gezielt eine einzelne Frequenz.

Anwendung:

\begin{itemize}
\item 50 Hz Netzbrummen
\item 60 Hz Netzbrummen
\item Maschinenvibrationen
\end{itemize}

Frequenzgang:

\begin{center}
\includegraphics[width=0.7\linewidth]{images/notch_response.png}
\end{center}

%%%%%%%%%%%%%%%%%%%%%%%%%%%%%%%%%%%%%%%%%%%%%%%%%%%%%%%%%%%%
\subsection{Single Notch Filter}

Die wichtigste Formel lautet:

\[
H(z)=
b
\frac{
1-2\cos(\omega_0)z^{-1}+z^{-2}
}
{
1-2b\cos(\omega_0)z^{-1}
+
(2b-1)z^{-2}
}
\]

\[
\omega_0
\]

ist die Notch-Frequenz.

%%%%%%%%%%%%%%%%%%%%%%%%%%%%%%%%%%%%%%%%%%%%%%%%%%%%%%%%%%%%
\subsection{Q-Faktor eines Notch Filters}

Die Breite der Kerbe wird durch den Q-Faktor beschrieben:

\[
Q=
\frac{\omega_0}
{\Delta\omega}
\]

mit

\[
\Delta\omega
\]

als 3-dB-Bandbreite.

Daraus:

\[
\Delta\omega=
\frac{\omega_0}{Q}
\]

\paragraph{Interpretation}

\begin{table}[h]
\centering
\begin{tabular}{c|c}
Q & Wirkung \\
\hline
klein & breite Kerbe \\
groß & schmale Kerbe \\
\end{tabular}
\end{table}

%%%%%%%%%%%%%%%%%%%%%%%%%%%%%%%%%%%%%%%%%%%%%%%%%%%%%%%%%%%%
\subsection{Signal Averaging}

Signal Averaging wird verwendet, wenn mehrere identische Messungen vorhanden sind.

Beispiel:

\[
x(n)=s(n)+v(n)
\]

Messung 1:

\[
s(n)+v_1(n)
\]

Messung 2:

\[
s(n)+v_2(n)
\]

usw.

%%%%%%%%%%%%%%%%%%%%%%%%%%%%%%%%%%%%%%%%%%%%%%%%%%%%%%%%%%%%
\subsection{Mittelwertbildung}

Bei N Messungen:

\[
\hat s(n)
=========

\frac1N
\sum_{k=1}^{N}
x_k(n)
\]

Da

\[
E[v(n)]=0
\]

verschwindet das Rauschen im Mittel.

%%%%%%%%%%%%%%%%%%%%%%%%%%%%%%%%%%%%%%%%%%%%%%%%%%%%%%%%%%%%
\subsection{Rauschreduktion durch Averaging}

Die Standardabweichung reduziert sich zu:

\[
\sigma_{\hat v}
===============

\frac{\sigma_v}
{\sqrt N}
\]

\paragraph{Beispiele}

\[
N=4
\]

\[
\sigma_{\hat v}
===============

\frac{\sigma_v}{2}
\]

\[
N=100
\]

\[
\sigma_{\hat v}
===============

\frac{\sigma_v}{10}
\]

%%%%%%%%%%%%%%%%%%%%%%%%%%%%%%%%%%%%%%%%%%%%%%%%%%%%%%%%%%%%
\subsection{Circular Buffer Averaging}

Für Aufgabe 8.42 wird ein Circular Buffer verwendet.

Aktualisierung:

\[
q_D=(q+D)\bmod(D+1)
\]

\[
w[q]
====

w[q_D]
+
x
\]

Ausgabe:

\[
y=
\frac{w[q]}{N}
\]

Nur ein Buffer-Eintrag wird pro Sample verändert.

%%%%%%%%%%%%%%%%%%%%%%%%%%%%%%%%%%%%%%%%%%%%%%%%%%%%%%%%%%%%
\subsection{Savitzky-Golay Filter}

Ein gewöhnlicher Moving-Average-Filter glättet zwar das Rauschen,
verzerrt jedoch Peaks.

Savitzky-Golay verwendet stattdessen lokale Polynomanpassung.

%%%%%%%%%%%%%%%%%%%%%%%%%%%%%%%%%%%%%%%%%%%%%%%%%%%%%%%%%%%%
\subsection{Grundidee}

Betrachte ein Fenster:

\[
x(n-2)
\quad
x(n-1)
\quad
x(n)
\quad
x(n+1)
\quad
x(n+2)
\]

Nun wird ein Polynom angepasst:

\[
p(t)
====

a+bt+ct^2
\]

Das Polynom approximiert die Datenpunkte im Least-Squares-Sinn.

%%%%%%%%%%%%%%%%%%%%%%%%%%%%%%%%%%%%%%%%%%%%%%%%%%%%%%%%%%%%
\subsection{Vorteile von Savitzky-Golay}

\begin{table}[h]
\centering
\begin{tabular}{c|c}
Moving Average & Savitzky-Golay \\
\hline
Peaks werden kleiner & Peaks bleiben erhalten \\
starke Glättung & moderate Glättung \\
verzerrt Krümmungen & erhält Krümmungen \\
\end{tabular}
\end{table}

%%%%%%%%%%%%%%%%%%%%%%%%%%%%%%%%%%%%%%%%%%%%%%%%%%%%%%%%%%%%
\subsection{Prüfungsrelevante Formelsammlung}

\[
x(n)=s(n)+v(n)
\]

\[
NRR=
\frac{\sigma_y^2}{\sigma_x^2}
=============================

\sum h^2(n)
\]

\[
SNR=
\frac{E[s^2]}{E[v^2]}
\]

\[
H(z)=1\pm z^{-D}
\]

\[
z_k=e^{j2\pi k/D}
\]

\[
H(z)=
\frac{b}
{1-az^{-1}}
\]

\[
NRR=
\frac{1-a}{1+a}
\]

\[
Q=
\frac{\omega_0}{\Delta\omega}
\]

\[
\sigma_{\hat v}
===============

\frac{\sigma_v}{\sqrt N}
\]

\[
\hat s(n)
=========

\frac1N
\sum x_k(n)
\]

\[
y=(1-\alpha)w[k]+\alpha w[k+1]
\]
